The following figure of merit (FOM) for nuclear material attractiveness is taken directly from a paper by Lawrence Livermore National Laboratory \cite{King2010TheScenario}. 

\begin{equation}
    \text{FOM}_1 = 1-\log\left( M\left[\frac{1}{800} + \frac{h}{4500}\right] +\frac{M}{50}\left[\frac{D}{500}\right]^{1/\log2}\right)
\end{equation}

\begin{displayquote}
Where M is the bare critical mass of the metal product material (impure plutonium) that is intended for the nuclear explosive in kg, h is the heat content in W/kg, and D is the dose rate of a 0.2M sphere evaluated at 1 m from the surface in rad/h. Table 1 correlates the result for the FOM calculation with the material's nuclear explosive utility (Bathke et al., 2009). In this FOM, the heat source h and the dose rate D are those that result directly from the isotopic composition of the metal in its transformed state after removal of fission products. The bare critical mass M is the critical mass after the same transformation. bare critical mass (kg) of the seperated element of interest \cite{King2010TheScenario}.
\end{displayquote}

\begin{table}[H]
\centering
\caption{Correlation between FOM$_1$ and nuclear explosives utility \cite{King2010TheScenario}}
\label{tab:figure_labels}
\renewcommand{\arraystretch}{1.4}
\begin{tabular}{c l}
\toprule
\textbf{FOM$_1$} & \textbf{Nuclear explosives utility} \\
\midrule
$>$2 & Preferred \\
1-2 & Usable \\
0-1 & Unsable \\
$<$0 & Unsable \\
\bottomrule
\end{tabular}
\end{table}